\documentclass[11pt]{article}
\usepackage[a4paper,margin=28mm]{geometry}
\usepackage{amsmath,amssymb,amsthm}
\usepackage[colorlinks=true,linkcolor=blue,urlcolor=blue,citecolor=blue]{hyperref}
\newtheorem*{theorem*}{Theorem}
\newtheorem*{corollary*}{Corollary}
\newenvironment{statement}[1]{\par\medskip\noindent\phantomsection\textbf{#1}\ }{\par\medskip}
\setlength{\emergencystretch}{3em}
\newcommand{\Q}{\mathbb Q}
\newcommand{\Z}{\mathbb Z}
\newcommand{\F}{\mathbb F}
\newcommand{\sd}{\mathbin{\rtimes}}
\DeclareMathOperator{\End}{End}
\DeclareMathOperator{\Gal}{Gal}
\DeclareMathOperator{\GL}{GL}
\DeclareMathOperator{\PGL}{PGL}
\DeclareMathOperator{\Hom}{Hom}
\DeclareMathOperator{\Res}{Res}
\DeclareMathOperator{\Norm}{N}
\DeclareMathOperator{\Ad}{Ad}
\DeclareMathOperator{\tr}{tr}
\DeclareMathOperator{\Tr}{Tr}
\title{Rational torsion containing \((\Z/3)^2\) forces \(2\)-primary torsion of order at most \(4\) on principally polarized abelian surfaces that are geometrically split and simple over \(\Q\)}
\author{}
\date{15 November 2026}

\begin{document}
\maketitle

\begin{theorem*}
Let \(A/\Q\) be an abelian surface with a principal polarization defined over \(\Q\), and suppose that \(A\) is
simple over \(\Q\) and geometrically split. If \((\Z/3)^2\subseteq A(\Q)\), then
\[\#A(\Q)[2^\infty]\le4 .\]
\end{theorem*}

\begin{corollary*}
No genus-\(2\) Jacobian over \(\Q\) that is geometrically split and simple over \(\Q\) has a rational torsion
subgroup that contains \((\Z/3)^2\) and at least eight points of \(2\)-power order. In particular
\href{https://genus-2-torsion.github.io/group.html?g=6.12}{\([6,12]\)},
\href{https://genus-2-torsion.github.io/group.html?g=3.24}{\([3,24]\)} and
\href{https://genus-2-torsion.github.io/group.html?g=2.6.6}{\([2,6,6]\)} do not occur as the rational torsion
subgroup of such a Jacobian.
\end{corollary*}

Here ``geometrically split'' means that \(A_{\bar\Q}\) is isogenous to a product of two elliptic curves; we impose
no restriction on the degree of this isogeny or on the field of definition of the geometric endomorphisms. The
Jacobian of a curve of genus \(2\) over \(\Q\) carries its canonical principal polarization, which is defined over
\(\Q\). The three groups of the Corollary have order \(72\) and \(2\)-primary parts of order \(8\); among the groups
of the census on 15 November 2026 they are the only ones that contain \((\Z/3)^2\) and have a \(2\)-primary part
of order at least \(8\). The theorem is the analogue for \((\Z/3)^2\) of the bound \(\#A(\Q)[2^\infty]\le4\) for surfaces with a
rational point of order \(5\) \cite{P40}. We do not know whether the bound \(4\) is attained: in the LMFDB, the
curves of this class whose Jacobians have \((\Z/3)^2\subseteq J(\Q)_{\mathrm{tors}}\) have at most two rational
points of \(2\)-power order (\hyperref[rem-lmfdb]{Remark~4}).

\emph{Outline.} We follow the case distinction of \cite{P63} and \cite{P2214} into types (N), (CM) and (S)
(Section~\ref{inputs}). For a polarization over \(\Q\) of degree prime to \(3\), a rational subgroup
\((\Z/3)^2\) is Lagrangian for the Weil pairing on \(A[3]\). So the \(3\)-division field is a generalized
dihedral extension \(\Q(\zeta_3,\sqrt[3]{q_1},\dots,\sqrt[3]{q_k})\), whose only quadratic subfield is
\(\Q(\sqrt{-3})\), and by Silverberg's theorem all geometric endomorphisms are defined over it
(\hyperref[thm-full3]{Theorem~1}). Hence the two elliptic curves of type (N) are defined over
\(\Q(\sqrt{-3})\), the CM field in type (CM) is \(\Q(\sqrt{-3})\), and in type (S) the Galois group \(\Gamma\) of
the endomorphism field \(L\) is cyclic of order \(2\) or dihedral of order \(6\). In type (N) the principal
polarization makes \(A\) a conjugate gluing of an elliptic curve \(E\) over \(K=\Q(\sqrt{-3})\) with \(E^\sigma\),
and \((\Z/3)^2\) forces \(E[3]\subseteq E(K)\); then \(E(K)_{\mathrm{tors}}\) is \(\Z/3\times\Z/3\) or
\(\Z/3\times\Z/6\), and \(\#A(\Q)[2^\infty]\le\#E(K)[2^\infty]^2\le4\) (\hyperref[thm-N]{Theorem~5}). In type
(CM), reduction at a prime \(p\equiv5,7\pmod{12}\) inert in the CM field gives a supersingular surface, whose
number of points is never divisible by \(72\) (\hyperref[thm-cm]{Theorem~4}). If \(\Gamma\) is dihedral of order
\(6\), a Frobenius element of order \(3\) at a prime \(p\equiv3\pmod4\) gives
\(8\nmid\#\widetilde A(\F_p)\) (\hyperref[lemma-shapes]{Lemma~2} and \hyperref[prop-ord3]{Proposition~3}). If
\(\Gamma\) has order \(2\), then \(L=\Q(\sqrt{-3})\) is imaginary, which forces real multiplication
(\hyperref[lemma-rm]{Lemma~6}). Then an element \(h\) of an inertia group at \(3\) satisfies
\(\rho_\ell(h)^2=-1\) for every \(\ell\neq3\), and \(A(\Q)[2^\infty]\) lies in \(A[2^\infty]^h\cong(\Z/2)^2\)
(\hyperref[thm-inertia]{Theorem~8}). The key step there is potentially good reduction at \(3\): a real
multiplication restricts to an isogeny \(E\to E^\sigma\) of non-square degree, which is impossible when
\(j(E)\) and \(j(E^\sigma)\) have the same negative valuation (\hyperref[lemma-tate]{Lemma~7}).

\section{Conventions and inputs}\label{inputs}

All number fields lie in a fixed algebraic closure \(\bar\Q\), and \(G_k=\Gal(\bar\Q/k)\). For a prime \(\ell\),
\(\chi_\ell\) is the \(\ell\)-adic cyclotomic character, \(T_\ell\) the Tate module and
\(V_\ell=T_\ell\otimes\Q_\ell\); \(\bar\chi_3\colon G_{\Q}\to\F_3^\times=\{\pm1\}\) is the mod-\(3\) cyclotomic
character, and \(\zeta\) is a primitive cube root of unity. For an abelian variety \(X\),
\(\End^0=\End\otimes\Q\), and \(g\in G_{\Q}\) acts on geometric endomorphisms by \(g(\varphi)=g\circ\varphi\circ
g^{-1}\). For a quadratic field \(K\), \(\sigma\) is its nontrivial automorphism and \(\tilde\sigma\in G_{\Q}\) a
lift; for a variety \(Y\) over \(K\), \(Y^\sigma\) is the conjugate variety and \(\tilde\sigma\colon
Y(\bar\Q)\to Y^\sigma(\bar\Q)\) the induced bijection \cite[Section~1]{P2214}. A polarization \(\lambda\colon
A\to A^\vee\) is defined over \(\Q\) if the isogeny \(\lambda\) is; every abelian variety over \(\Q\) has a
polarization defined over \(\Q\). \(\Q^{\times2}\) is the set of nonzero rational squares, and \(v_\ell\) the
\(\ell\)-adic valuation.

\begin{itemize}
\item \textbf{(I1) Torsion over quadratic fields.} By Kenku--Momose \cite[Theorem~(0.1)]{KM} and Kamienny
\cite{Ka} (see the statement in \cite[Section~1]{Na}), a torsion subgroup of an elliptic curve over a quadratic
field that contains \(\Z/3\times\Z/3\) is isomorphic to \(\Z/3\times\Z/3\) or \(\Z/3\times\Z/6\). By the Weil
pairing this happens only over \(\Q(\sqrt{-3})\), where \cite[Theorem~2(ii)]{Na} lists the same two groups.

\item \textbf{(I2) Reduction} \cite{ST}, as stated in \cite[Section~1]{P63} and \cite[(I2)]{P2214}: let \(K'\) be
a finite extension of \(\Q_p\) with residue field \(k\), \(X\) an abelian variety over \(K'\), and \(\ell\neq p\)
a prime. (a) Good reduction is detected by the inertia action on \(T_\ell(X)\) \cite[\S1, Theorem~1]{ST}; it is
preserved by \(K'\)-isogenies \cite[\S1, Corollary~2]{ST}. (b) If \(X\) has good reduction \(\widetilde X\) and
\(m\) is prime to \(p\), then \(X(K')[m]\) injects into \(\widetilde X(k)\) \cite[Lemma~2]{ST}. (c) If \(X\) has
potentially good reduction, then it has good reduction over a finite extension of \(K'\) with residue field
\(k\) \cite[\S2]{ST}. (d) If \(X\) has potentially good reduction and \(g\) lies in the inertia group of \(K'\),
then \(\rho_\ell(g)\) has finite order, and its characteristic polynomial on \(V_\ell(X)\) has integer coefficients
that do not depend on \(\ell\) (apply \cite[\S2, Theorem~2]{ST} to the powers of \(g\)). (e) (Weil, Tate.) For an abelian surface \(B\) over
\(\F_q\), \(\#B(\F_q)=P(1)\), where \(P\in\Z[X]\) is the characteristic polynomial of the Frobenius endomorphism
\(\pi\) on \(V_\ell(B)\); the complex roots of \(P\) have absolute value \(\sqrt q\); and \(\Q[\pi]\) is the centre
of the semisimple algebra \(\End^0_{\F_q}(B)\) \cite[Theorem~2(a)]{Ta}, hence a product of fields. By (b) and (e),
if an abelian surface \(X\) over a number field \(k\) has good reduction \(\widetilde X\) at a prime
\(\mathfrak p\nmid\ell\), then \(\#\widetilde X(O_k/\mathfrak p)=\det(1-\mathrm{Frob}_{\mathfrak p}\mid V_\ell)\)
for an arithmetic Frobenius element \(\mathrm{Frob}_{\mathfrak p}\in G_k\) at \(\mathfrak p\)
\cite[(I1)]{P2210}.

\item \textbf{(I3) The case distinction} \cite[(I3)]{P2214}, from \cite[Sections~2--5]{P63}. A geometrically
split abelian surface \(A/\Q\) is of exactly one of the types \textbf{(N)} (\(A_{\bar\Q}\) is isogenous to
\(E_1\times E_2\) with \(E_1\), \(E_2\) not isogenous), \textbf{(CM)} (\(A_{\bar\Q}\) is isogenous to \(E^2\) with
\(E\) a CM curve) and \textbf{(S)} (\(A_{\bar\Q}\) is isogenous to \(E^2\) with \(E\) without CM). In type (S),
\(D=\End^0(A_{\bar\Q})\cong M_2(\Q)\); \(L\) is the smallest field over which all geometric endomorphisms are
defined, a finite Galois extension of \(\Q\); \(\Gamma=\Gal(L/\Q)\) acts faithfully on \(D\), the element \(g\)
acting by \(\Ad(R_g)\colon\varphi\mapsto R_g\varphi R_g^{-1}\) with \(R_g\in D^\times\) unique up to \(\Q^\times\);
so \(\Gamma\) embeds in \(\PGL_2(\Q)\), and \(\End^0_{\Q}(A)=D^\Gamma\). The group \(\Gamma\) is one of \(C_1\),
\(C_2\), \(C_3\), \(C_4\), \(C_6\), \(V_4\), \(D_6\), \(D_8\), \(D_{12}\) (\(D_{2m}\) is dihedral of order \(2m\)),
and \(\Gamma\neq C_1\) if \(A\) is simple over \(\Q\). An element of \(\Gamma\) of order \(n\in\{3,4,6\}\) has a
unique representative \(R\) with \(\tr R=1\), and \(R^2-R+1/c=0\) with \(c=1,2,3\) respectively; an element of
order \(2\) has a representative \(R\) with \(\tr R=0\), and then \(R^2\in\Q^\times\) \cite[Section~5]{P63}. In
type (S), every elliptic curve in \(A_{\bar\Q}\) is defined over \(L\), and any two of them are \(L\)-isogenous
\cite[Lemma~3]{P63}. In type (N), if \(A\) is simple over \(\Q\), then \(A_{\bar\Q}\) contains exactly two
elliptic curves; they are defined over a quadratic field \(K\) and exchanged by \(\tilde\sigma\)
\cite[Section~3]{P63}.

\item \textbf{(I4) The determinant} \cite[Lemma~3]{P2214}, in the form \cite[(I6)]{P2210}. Let \(A\) be of type
(S) with \(\Gamma=\langle r\rangle\) cyclic of order \(n\in\{3,4,6\}\), and \(R\) the trace-one representative of
\(r\). Then \(F=\Q[R]\) is a field (\(F\cong\Q(\sqrt{-3})\) for \(n=3,6\), and \(F\cong\Q(i)\) for \(n=4\)),
\(F=D^\Gamma=\End^0_{\Q}(A)\), and \(\Norm(R)=\det R=1/c>0\). For an elliptic curve \(E\subset A_L\) and
\(H=\Hom^0(E_{\bar\Q},A_{\bar\Q})\), a one-dimensional \(F\)-vector space, evaluation is an \(F\)-linear isomorphism
\(V_\ell(E)\otimes_{\Q}H\to V_\ell(A)\); so \(V_\ell(A)\) is free of rank \(2\) over \(F\otimes\Q_\ell\), on which
\(G_{\Q}\) acts \(F\otimes\Q_\ell\)-linearly. Under this isomorphism each \(g\in G_{\Q}\) acts as
\(B_g\otimes C_g\), where \(B_g\in\GL(V_\ell(E))\) and \(C_g\in\GL(H)=D^\times\) satisfies
\(g(\varphi)=C_g\varphi C_g^{-1}\) for \(\varphi\in D\); thus \(C_g\) is multiplication by an element
\(\gamma_g\in F^\times\). Moreover \(\varepsilon(g)=\gamma_g^2/\deg(u_g)\), for a quasi-isogeny \(u_g\colon E\to
E^g\) of degree \(\deg(u_g)\in\Q_{>0}\), defines a character \(\varepsilon\colon G_{\Q}\to F^\times\) of finite
order that factors through \(\Gamma\), with \(\varepsilon(r)=R/\bar R\) a primitive \(n\)-th root of unity; and the
\(F\otimes\Q_\ell\)-determinant of \(g\) is \(\chi_\ell(g)\varepsilon(g)\). The proof uses only these properties of
\(F\) and \(R\), and that the roots of unity in \(F\) have norm \(1\).

\item \textbf{(I5) Weil restriction} \cite[(I7)]{P45}. For an elliptic curve \(E\) over a quadratic field \(K\),
\(\Res_{K/\Q}E\) is an abelian surface over \(\Q\) whose \(\bar\Q\)-points are the pairs \((x,y)\) with
\(x\in E(\bar\Q)\) and \(y\in E^\sigma(\bar\Q)\); an element \(g\in G_K\) maps \((x,y)\) to \((gx,gy)\), an element
\(g\notin G_K\) maps it to \((gy,gx)\), and the rational points are the pairs \((x,\sigma x)\) with
\(x\in E(K)\).

\item \textbf{(I6) Conjugate gluing} \cite[proof of Lemma~3]{P45}, the paragraphs ``Each is the quotient by the
other'' and ``The gluing'', which use only a principal polarization \(\lambda\) of \(A\) defined over \(\Q\) and
the two elliptic curves \(E\), \(E^\sigma\) of type (N) (not the point of order \(45\)). The complement of \(E\)
for \(\lambda\) is \(E^\sigma\); \(E\cap E^\sigma=E[n]=E^\sigma[n]\) for some \(n\ge1\); the map
\(\pi_E=(\lambda^0_E)^{-1}i_E^\vee\lambda\colon A_K\to E\) (with \(i_E\) the inclusion and \(\lambda^0_E\) the
canonical principal polarization of \(E\)) is surjective with kernel \(E^\sigma\), so that
\(A_K/E^\sigma\cong E\); the addition map \(\varphi_n\colon E\times E^\sigma\to A_K\) has as kernel the graph of
the anti-isometry \(\psi_n\colon E[n]\to E^\sigma[n]\), \(x\mapsto-x\); \(\Phi_n=(\pi_E,\pi_{E^\sigma})\) satisfies
\(\Phi_n\varphi_n=[n]\) and \(\varphi_n\Phi_n=[n]\); and \(\varphi_n\) and \(\Phi_n\) descend to isogenies
\(\Res_{K/\Q}E\to A\) and \(A\to\Res_{K/\Q}E\) over \(\Q\), with the Galois action of (I5).

\item \textbf{(I7) Deuring's criterion} \cite[Chapter~13, Theorem~12]{La}, due to Deuring \cite{De}: an elliptic
curve over a number field with complex multiplication by an order in an imaginary quadratic field \(M\), with good
reduction at a prime above a prime \(p\) that is inert in \(M\), has supersingular reduction there.

\item \textbf{(I8) The \(j\)-function at a prime} \cite[Theorem~V.3.1(b) and Lemma~V.5.1]{Si2}, and potentially
good reduction \cite[Proposition~VII.5.5]{Si1}. Let \(j=q^{-1}+\sum_{n\ge0}c(n)q^n\), \(c(n)\in\Z\), be the
\(q\)-expansion of the elliptic modular function. Let \(K'\) be a finite extension of \(\Q_p\) and \(v\) the
valuation of \(\bar K'\) with \(v(p)=1\). For \(q\) in a finite extension of \(K'\) with \(v(q)>0\), the series
\(j(q)=q^{-1}+\sum c(n)q^n\) converges, and \(v(j(q))=-v(q)\). Every \(\alpha\in\bar K'\) with \(v(\alpha)<0\)
equals \(j(q)\) for a unique \(q\in\Q_p(\alpha)\) with \(v(q)>0\). An elliptic curve \(E\) over \(K'\) has
potentially good reduction if and only if \(v(j(E))\ge0\).

\item \textbf{(I9) The modular polynomial} \cite[Chapter~5, Sections~1--3]{La}. For \(N\ge1\) the modular
polynomial \(\Phi_N\in\Z[X,Y]\) satisfies, in \(\Q(\zeta_N)((q^{1/N}))[X]\),
\[\Phi_N(X,j(q))=\prod_{s,b,d}\bigl(X-j(\zeta_d^b\,q^{s/d})\bigr),\]
the product over the integers \(s,d>0\) with \(sd=N\) and \(0\le b<d\) with \(\gcd(s,b,d)=1\), where
\(\zeta_d=\zeta_N^{N/d}\) and \(j(\zeta_d^bq^{s/d})\) is the series of (I8) evaluated at \(\zeta_d^bq^{s/d}\)
(formula (1) of \cite[Chapter~5, Section~2]{La}). If \(B\to A\) is a cyclic isogeny of degree \(N\) of elliptic
curves over a field of characteristic \(0\), then \(\Phi_N(j(B),j(A))=0\) \cite[Chapter~5, Theorem~5 and the
remark following it]{La}.
\end{itemize}

\section{A rational subgroup of order nine in the \(3\)-torsion}\label{full3}

\begin{statement}{Theorem 1 (rational 3-torsion of rank two).}\label{thm-full3}
Let \(A/\Q\) be an abelian surface with a polarization defined over \(\Q\) of degree prime to \(3\) (for example a
principal polarization defined over \(\Q\)), and suppose that \(A(\Q)\) contains a subgroup \(W\cong(\Z/3)^2\).
\begin{enumerate}
\item[(a)] \(W\) is a Lagrangian subgroup of \(A[3]\), and in a symplectic basis of \(A[3]\) whose first two
vectors span \(W\) every \(g\in G_{\Q}\) acts by
\[\begin{pmatrix} I & X_g\\ 0&\bar\chi_3(g)I\end{pmatrix}\qquad\text{with }X_g\text{ symmetric.}\]
\item[(b)] \(\Q(A[3])=\Q(\zeta)(\sqrt[3]{q_1},\dots,\sqrt[3]{q_k})\) with \(q_i\in\Q^\times\) independent in
\(\Q^\times/\Q^{\times3}\), and
\(\Gal(\Q(A[3])/\Q)\cong U\sd C_2\) with \(U\) elementary abelian of order \(3^k\le27\), on which \(C_2\) acts by
inversion. Every quotient of \(U\sd C_2\) is trivial or again of this form; in particular the only quadratic
subfield of \(\Q(A[3])\) is \(\Q(\sqrt{-3})\).
\item[(c)] Every endomorphism and every abelian subvariety of \(A_{\bar\Q}\) is defined over \(\Q(A[3])\).
\item[(d)] If \(A\) is simple over \(\Q\) and geometrically split, then: in type (N), \(K=\Q(\sqrt{-3})\); in type
(CM), the CM field is \(\Q(\sqrt{-3})\); in type (S), either \(\Gamma\cong C_2\), \(L=\Q(\sqrt{-3})\) and
\(\End^0_{\Q}(A)\cong\Q(\sqrt a)\) with \(a>0\) not a square (so the generator of \(\Gamma\) acts by
\(\Ad(R)\) with \(R^2>0\)), or \(\Gamma\cong D_6\) and the unique quadratic subfield of \(L\) is
\(\Q(\sqrt{-3})\).
\end{enumerate}
\end{statement}

\emph{Proof.} (a) Let \(e\) be the Weil pairing on \(A[3]\) attached to the polarization; it is nondegenerate
because the degree is prime to \(3\), alternating, and \(e(gx,gy)=g(e(x,y))=e(x,y)^{\bar\chi_3(g)}\). For
\(x,y\in W\) we get \(e(x,y)\in\mu_3^{G_{\Q}}=1\), so \(W\) is isotropic of dimension \(2\), hence Lagrangian, and
\(x\mapsto e(x,\cdot)|_W\) identifies \(A[3]/W\) with \(\Hom(W,\mu_3)\cong\mu_3^2\) as Galois modules. Choose a basis
\(w_1,w_2\) of \(W\) and complete it by \(w_1',w_2'\) with \(e(w_i,w_j')=\zeta^{\delta_{ij}}\) and
\(e(w_1',w_2')=1\) (a symplectic basis). Then \(g\) acts trivially on \(W\) and by \(\bar\chi_3(g)\) on \(A[3]/W\),
which gives the displayed shape; writing \(e\) additively through \(\zeta\mapsto1\), the matrix \(g\) multiplies the
Gram matrix \(\left(\begin{smallmatrix}0&I\\-I&0\end{smallmatrix}\right)\) by \(\bar\chi_3(g)\), which for
matrices of this shape says that \(X_g\) is symmetric (Part~A1 of the certificate).

(b) Put \(k_3=\Q(\zeta)\subseteq\Q(A[3])\). By (a) the image \(G\) of \(G_{\Q}\) in \(\GL(A[3])\) is contained in the
group \(M\) of matrices \(\left(\begin{smallmatrix}I&X\\0&\pm I\end{smallmatrix}\right)\) with \(X\) symmetric, of
order \(54\), and \(U=G\cap\{\left(\begin{smallmatrix}I&X\\0&I\end{smallmatrix}\right)\}\) is the image of
\(G_{k_3}\); it is elementary abelian of order at most \(27\), and of index \(2\) in \(G\) because
\(k_3\neq\Q\). Every element \(t\in G\setminus U\) has the form
\(\left(\begin{smallmatrix}I&Y\\0&-I\end{smallmatrix}\right)\) and satisfies \(t^2=1\) and \(tut^{-1}=u^{-1}\) for
\(u\in U\), so \(G=U\sd\langle t\rangle\) with \(t\) acting by inversion. By Kummer theory
\(\Q(A[3])=k_3(\sqrt[3]{\beta}:\beta\in B)\) for a subgroup \(B\subset k_3^\times/k_3^{\times3}\); equivariance of
the Kummer pairing under \(t\), which inverts both \(U\) and \(\mu_3\), shows that \(B\) is fixed by complex
conjugation. If \(\bar\beta=\beta\gamma^3\), then \(\Norm(\beta)=\beta^2\gamma^3\) and
\(\beta\equiv\Norm(\beta)^{-1}\pmod{k_3^{\times3}}\), with \(\Norm(\beta)\in\Q^\times\). Quotients: if
\(N\trianglelefteq G\) contains some \(tu\), then for \(v\in U\) it contains \((tu)^{-1}v(tu)v^{-1}=v^{-2}=v\), so
\(N\supseteq U\), \(t\in N\) and \(N=G\). Otherwise \(N\subseteq U\) and \(G/N\cong(U/N)\sd C_2\). The only subgroup
of index \(2\) of \(G\) is \(U\), whose fixed field is \(k_3=\Q(\sqrt{-3})\). Part~A of the certificate checks these
statements for every subgroup \(G\subseteq M\) that contains an element with \(-I\) in the lower right corner
(A2), and for the abstract groups \(U\sd C_2\) with \(\#U\le27\) (A3).

(c) (Silverberg's theorem \cite[Theorem~2.4]{Sb}; we include the short proof.) Let \(s\in G_{\Q(A[3])}\). For
\(\varphi\in\End(A_{\bar\Q})\) and \(x\in A[3]\), \(s(\varphi)(x)=s(\varphi(s^{-1}x))=\varphi(x)\), so
\(s(\varphi)-\varphi\) kills \(A[3]\) and lies in \(3\End(A_{\bar\Q})\). Thus \(s\) acts on the free \(\Z\)-module
\(\End(A_{\bar\Q})\) by an automorphism of finite order congruent to \(1\) modulo \(3\), which is the identity
(Minkowski). Every abelian subvariety is the image of an endomorphism (Poincar\'e).

(d) In type (N) the two elliptic curves are defined over \(\Q(A[3])\) by (c), so \(K\subseteq\Q(A[3])\), and
\(K=\Q(\sqrt{-3})\) by (b). In type (CM), \(A_{\bar\Q}\) is isogenous to \(E^2\) with \(E\) a CM curve, and
\(\End^0(A_{\bar\Q})\cong M_2(M)\) for the CM field \(M\) of \(E\). Let \(L\) be the field of definition of
\(\End(A_{\bar\Q})\); then \(L\subseteq\Q(A[3])\) by (c). The centre \(M\) of \(\End^0(A_{\bar\Q})\) acts on the
two-dimensional tangent space of \(A_L\) at the origin, an \(L\)-vector space, through one embedding \(\iota\) of
\(M\) (with multiplicity two), and the characteristic polynomial \((X-\iota(\mu))^2\) of an element \(\mu\in M\)
has coefficients in \(L\). So \(\iota(M)\subseteq L\subseteq\Q(A[3])\), and \(M\cong\Q(\sqrt{-3})\) by (b). In
type (S), \(L\subseteq\Q(A[3])\) by (c), so \(\Gamma\) is a quotient of \(G\). By (b), in the list of (I3) only
\(C_1\), \(C_2\) and \(D_6\) are trivial or of the form \(U\sd C_2\) (Part~A of the certificate), and \(C_1\) is
excluded by simplicity. If \(\Gamma\cong C_2\), then \(L\) is the quadratic subfield \(\Q(\sqrt{-3})\) of
\(\Q(A[3])\). It is imaginary, so by \hyperref[lemma-rm]{Lemma~6} below (the argument of
\cite[Section~9.1]{P63}) the generator of \(\Gamma\) acts by \(\Ad(R)\) with \(R^2=a>0\) not a square, and
\(\End^0_{\Q}(A)=D^\Gamma=\Q(R)\cong\Q(\sqrt a)\). If \(\Gamma\cong D_6\), the image of \(U\) in \(\Gamma\) is a
normal \(3\)-subgroup, and \(G/U\to\Gamma/\mathrm{im}(U)\) is onto, so \(\mathrm{im}(U)\cong C_3\) and
\(L^{\mathrm{im}(U)}\subseteq\Q(A[3])^U=\Q(\sqrt{-3})\); the unique quadratic subfield of \(L\) is therefore
\(\Q(\sqrt{-3})\). \qed

\section{Frobenius elements of order \(3\), \(4\) and \(6\)}\label{frobenius}

\begin{statement}{Lemma 2 (Frobenius polynomials).}\label{lemma-shapes}
Let \(A/\Q\) be of type (S), \(\gamma\in\Gamma\) of order \(n\in\{3,4,6\}\), and \(p\) a prime of good reduction
for \(A\), unramified in \(L\), whose Frobenius class in \(\Gamma\) contains \(\gamma\). Then the characteristic
polynomial \(P\) of Frobenius of \(\widetilde A/\F_p\) and \(\#\widetilde A(\F_p)=P(1)\) are, for a suitable choice
of the primitive cube root of unity \(\zeta\) and integers \(a\), \(b\), \(u\):
\[\begin{array}{lll}
n=3\colon & X^4+aX^3+(a^2-p)X^2+apX+p^2, & \#\widetilde A(\F_p)=\Norm\bigl((1+a)+(a+p)\zeta\bigr),\\[2pt]
n=4\colon & X^4-2uX^3+2u^2X^2-2puX+p^2, & \#\widetilde A(\F_p)=(1-u)^2+(p-u)^2,\\[2pt]
n=6\colon & X^4+3bX^3+(3b^2+p)X^2+3bpX+p^2, & \#\widetilde A(\F_p)=\Norm\bigl((1+2b+p)+(b+p)\zeta\bigr),
\end{array}\]
with \(\Norm=\Norm_{\Q(\zeta)/\Q}\). Consequently:
\begin{enumerate}
\item[(i)] if \(p\equiv2\pmod3\), then \(v_3(\#\widetilde A(\F_p))\le1\) for \(n=3\), and
\(3\nmid\#\widetilde A(\F_p)\) for \(n\in\{4,6\}\);
\item[(ii)] if \(p\equiv3\pmod4\), then \(v_2(\#\widetilde A(\F_p))\in\{0,2\}\) for \(n\in\{3,6\}\), and
\(8\nmid\#\widetilde A(\F_p)\) for \(n=4\).
\end{enumerate}
\end{statement}

\emph{Proof.} Let \(L'=L^{\langle\gamma\rangle}\), \(\mathfrak P\mid p\) a prime of \(L\) with Frobenius
\(\gamma\), and \(\mathfrak p=\mathfrak P\cap L'\). The decomposition field of \(\mathfrak P\) is \(L'\), so
\(\mathfrak p\) has residue field \(\F_p\), and \(g=\mathrm{Frob}_{\mathfrak p}\in G_{L'}\) satisfies
\(\#\widetilde A(\F_p)=\det(1-g\mid V_\ell(A))\) by (I2)(e), for every prime \(\ell\neq p\). Put
\(F=\Q[R_\gamma]\subseteq D^{\langle\gamma\rangle}=\End^0_{L'}(A)\), where \(R_\gamma\) is the trace-one
representative of \(\gamma\); \(F\cong\Q(\sqrt{-3})\) for \(n=3,6\), and \(F\cong\Q(i)\) for \(n=4\). The proof of
\cite[Lemma~3]{P2214}, that is (I4), applies verbatim over \(L'\), with the cyclic group
\(\langle\gamma\rangle=\Gal(L/L')\) in place of \(\Gamma\): for \(g'\in G_{L'}\) the element \(C_{g'}\) represents a
power of \(\Ad(R_\gamma)\), so \(C_{g'}\in F\), and \(\Norm(R_\gamma)=1/c>0\). Hence \(V_\ell(A)\) is free of rank
\(2\) over \(F\otimes\Q_\ell\), the element \(g\) acts as \(B_g\otimes\gamma_g\) with \(\gamma_g\in F^\times\) and
\(\gamma_g^2=\varepsilon(g)\deg(u_g)\), \(\deg(u_g)\in\Q_{>0}\), and the \(F\otimes\Q_\ell\)-determinant of \(g\) is
\(p\,\varepsilon(g)\), where \(\varepsilon(g)=R_\gamma/\bar R_\gamma\) is a primitive \(n\)-th root of unity. Its
\(F\otimes\Q_\ell\)-trace is \(t=\tr(B_g)\gamma_g\). We choose \(\zeta\in F\) with \(\varepsilon(g)=\zeta\) if
\(n=3\) and \(\varepsilon(g)=-\zeta^2\) if \(n=6\), and \(i\in F\) with \(\varepsilon(g)=i\) if \(n=4\). Since
\(\gamma_g^2\in\varepsilon(g)\Q_{>0}\) is a square in \(F\), we get \(\gamma_g\in\Q\zeta^2\) (\(n=3\)),
\(\gamma_g\in\Q\zeta\sqrt{-3}\) (\(n=6\), where \(\zeta\sqrt{-3}=-2-\zeta\)) and \(\gamma_g\in\Q(1+i)\) (\(n=4\);
this is the computation of the proof of \cite[Proposition~A]{P2210}). Thus \(t=a\zeta^2\), \(t=b(-2-\zeta)\) and
\(t=u(1+i)\) with \(a,b,u\in\Q_\ell\), and
\[P(X)=\Norm_{F\otimes\Q_\ell/\Q_\ell}\bigl(X^2-tX+p\,\varepsilon(g)\bigr)\]
is the displayed polynomial (Part~B of the certificate, which also treats the other choices of \(\zeta\) and
\(i\)). The coefficients of \(P\) are integers by (I2)(e). For \(n=3\) the coefficient of \(X^3\) gives
\(a\in\Z\). For \(n=6\), \(3b\in\Z\) and \(3b^2+p\in\Z\) give \(3b^2=(3b)^2/3\in\Z\), so \(3\mid3b\) and
\(b\in\Z\). For \(n=4\), \(2u\in\Z\) and \(2u^2=(2u)^2/2\in\Z\) give \(u\in\Z\). Finally
\(\#\widetilde A(\F_p)=P(1)=\Norm(1-t+p\,\varepsilon(g))\), which is the displayed value. Statements (i) and (ii)
depend only on the residues of \((a,p)\), \((b,p)\), \((u,p)\) modulo \(9\), respectively \(16\), and are checked in
Part~C of the certificate (by hand: \(v_3(\Norm(x+y\zeta))\ge2\) if and only if \(3\mid x\) and \(3\mid y\), and
\(2\) is inert in \(\Q(\zeta)\)). \qed

\begin{statement}{Proposition 3.}\label{prop-ord3}
Let \(A/\Q\) be of type (S).
\begin{enumerate}
\item[(a)] If \(\Gamma\) has an element of order \(3\) and \(\sqrt{-3}\notin L\), then \(\#A(\Q)[3^\infty]\le3\).
\item[(b)] If \(\Gamma\in\{C_4,C_6\}\), then \(A(\Q)[3]=0\).
\item[(c)] If \(\Gamma\in\{C_3,C_6,D_6,D_{12}\}\) and \(\#A(\Q)[3^\infty]\ge9\), then \(\Gamma\in\{D_6,D_{12}\}\),
\(\sqrt{-3}\in L\), and \(\#A(\Q)[2^\infty]\le4\).
\item[(d)] If \(\Gamma=D_8\) and \(A(\Q)[3]\neq0\), then \(\Q(\sqrt{-3})=L^{\langle r\rangle}\) for a rotation \(r\)
of order \(4\), and \(\#A(\Q)[2^\infty]\le4\).
\end{enumerate}
\end{statement}

\emph{Proof.} In each case below, Chebotarev's theorem gives infinitely many primes with the stated properties,
so we may take them of good reduction and unramified in \(L\); for such a prime \(p\ne3\), \(A(\Q)[3^\infty]\)
and \(A(\Q)[2^\infty]\) (if \(p\) is odd) inject into \(\widetilde A(\F_p)\) by (I2)(b).

(a) As \(L\cap\Q(\zeta)=\Q\), Chebotarev's theorem in \(L(\zeta)\) gives infinitely many primes
\(p\equiv2\pmod3\) whose Frobenius class in \(\Gamma\) consists of elements of order \(3\); then
\(v_3(\#\widetilde A(\F_p))\le1\) by \hyperref[lemma-shapes]{Lemma~2}(i).

(b) \(C_6\): this is \cite[Lemma~9]{P63}. Directly: \(L\) is totally real \cite[Section~9.1]{P63}, so
\(\sqrt{-3}\notin L\), and Chebotarev's theorem in \(L(\zeta)\) gives primes \(p\equiv2\pmod3\) whose Frobenius
generates \(\Gamma\); \hyperref[lemma-shapes]{Lemma~2}(i) with \(n=6\) gives \(3\nmid\#\widetilde A(\F_p)\).
\(C_4\): a cyclic quartic field \(L\) does not contain \(\Q(\sqrt{-3})\), since complex conjugation would restrict
to an element of \(\Gal(L/\Q)\cong C_4\) acting nontrivially on the unique quadratic subfield, hence of order
\(4\), whereas complex conjugation has order at most \(2\) (as in \cite[proof of Proposition~A, case~(v)]{P45}).
So Chebotarev's theorem in \(L(\zeta)\) gives primes \(p\equiv2\pmod3\) whose Frobenius generates \(\Gamma\), and
\hyperref[lemma-shapes]{Lemma~2}(i) with \(n=4\) gives \(3\nmid\#\widetilde A(\F_p)\). (This case is also
\cite[proof of Proposition~A, case~(v)]{P45}.)

(c) By (a), \(\sqrt{-3}\in L\); \(C_3\) is impossible, as a cyclic cubic field has no quadratic subfield, and
\(C_6\) by (b). \(D_6\): \(D_6\) has a unique subgroup of index \(2\), so the only quadratic subfield of \(L\) is
\(\Q(\sqrt{-3})\) and \(i\notin L\); Chebotarev's theorem in \(L(i)\) gives primes \(p\equiv3\pmod4\) whose
Frobenius has order \(3\), and \hyperref[lemma-shapes]{Lemma~2}(ii) gives \(v_2(\#\widetilde A(\F_p))\le2\).
\(D_{12}\): let \(r\) generate the rotations (of order \(6\)). If \(r\) moved \(\sqrt{-3}\), the primes with
Frobenius \(r\) would be \(\equiv2\pmod3\), and \(3\nmid\#\widetilde A(\F_p)\) by
\hyperref[lemma-shapes]{Lemma~2}(i); so \(L^{\langle r\rangle}=\Q(\sqrt{-3})\). If \(i\notin L\), we use primes
\(p\equiv3\pmod4\) with Frobenius of order \(3\) (Chebotarev's theorem in \(L(i)\)); if \(i\in L\), then
\(i\notin L^{\langle r\rangle}\), so every prime with Frobenius \(r\) is \(\equiv3\pmod4\). In both cases
\hyperref[lemma-shapes]{Lemma~2}(ii) applies.

(d) Let \(r\) be a rotation of order \(4\). If \(\sqrt{-3}\notin L^{\langle r\rangle}\), there are infinitely many
primes \(p\equiv2\pmod3\) whose Frobenius class contains \(r\): by Chebotarev's theorem in \(L(\zeta)\) if
\(\sqrt{-3}\notin L\), and because every prime with Frobenius \(r\) is inert in \(\Q(\sqrt{-3})\) if
\(\sqrt{-3}\in L\). For them \hyperref[lemma-shapes]{Lemma~2}(i) with \(n=4\) gives \(3\nmid\#\widetilde A(\F_p)\),
contradicting \(A(\Q)[3]\neq0\). So \(\Q(\sqrt{-3})\subseteq L^{\langle r\rangle}\), which is quadratic, and
\(\Q(\sqrt{-3})=L^{\langle r\rangle}\). In the same way there are infinitely many primes \(p\equiv3\pmod4\) whose
Frobenius class contains \(r\): by Chebotarev's theorem in \(L(i)\) if \(i\notin L\); and if \(i\in L\), then
\(i\notin L^{\langle r\rangle}=\Q(\sqrt{-3})\), so every prime with Frobenius \(r\) is inert in \(\Q(i)\). For them
\hyperref[lemma-shapes]{Lemma~2}(ii) with \(n=4\) gives \(8\nmid\#\widetilde A(\F_p)\). \qed

\medskip
Only part (c), for \(\Gamma\cong D_6\), is used in the proof of the Theorem.

\section{The CM case}\label{cm}

\begin{statement}{Theorem 4 (the CM case).}\label{thm-cm}
Let \(A/\Q\) be an abelian surface such that \(A_{\bar\Q}\) is isogenous to \(E^2\) for an elliptic curve \(E\) with
complex multiplication. Then \(A(\Q)\) has no point of order \(9\) and no point of order \(12\), and
\(72\nmid\#A(\Q)_{\mathrm{tors}}\). No polarization and no simplicity over \(\Q\) is assumed.
\end{statement}

\emph{Proof.} Let \(M\) be the imaginary quadratic field such that \(E\) has complex multiplication by an order in
\(M\), and let \(p\ge5\) be a prime inert in \(M\). The \(j\)-invariant of \(E\) is an algebraic integer
\cite[Theorem~II.6.1]{Si2}, so \(E\) has potentially good reduction at every prime (I8). Over a finite extension
of \(\Q_p\), \(A\) becomes isogenous to \(E^2\) and \(E\) acquires good reduction; by (I2)(a), \(A\) acquires good reduction there, and by (I2)(c) it has
good reduction \(\widetilde A\) over a finite extension \(K'\) of \(\Q_p\) with residue field \(\F_p\). Over a
further finite extension, \(A\) is isogenous to \(E^2\) and \(E\) has good reduction \(\widetilde E\), which is
supersingular by (I7); reductions of isogenous abelian varieties with good reduction are isogenous, so
\(\widetilde A\) is isogenous over \(\bar\F_p\) to \(\widetilde E^2\). Over a finite field \(\F_{p^{2r}}\) over which
this isogeny and all endomorphisms of \(\widetilde E\) are defined, the Frobenius of \(\widetilde E\) is
\(\pm p^r\), so every Frobenius eigenvalue \(\alpha\) of \(\widetilde A\) over \(\F_p\) satisfies
\(\alpha^{2r}=\pm p^r\), and \(\alpha/\sqrt p\) is a root of unity. Write
\(P=X^4+c_1X^3+c_2X^2+pc_1X+p^2\) for the characteristic polynomial of Frobenius of \(\widetilde A\). Then
\(c_1/\sqrt p\) and \(c_2/p\) are algebraic integers, being sums of products of roots of unity, and \(c_1^2/p\) and
\(c_2/p\) are rational; so \(p\mid c_1\) and \(p\mid c_2\). Hence \(P=X^4+kpX^2+p^2\) with \(|k|\le2\) if
\(p\ge7\), and for \(p=5\) also \(P=X^4\pm5X^3+15X^2\pm25X+25\) (Part~D1 of the certificate for \(p\le23\); for
\(p\ge17\), \(p\mid c_1\) and \(|c_1|\le4\sqrt p<p\) give \(c_1=0\), and then \(|c_2|\le2p\) because the numbers
\(\alpha+p/\alpha=\pm\sqrt{2p-c_2}\) lie in \([-2\sqrt p,2\sqrt p]\)). The possible values of \(P(1)\) are
\((p+1)^2\), \(p^2+p+1\), \(p^2+1\), \(p^2-p+1\), \((p-1)^2\), and also \(11\) and \(71\) for \(p=5\) (Part~D2). The
values \(p^2\pm p+1\) are odd and never divisible by \(9\), and \(3\nmid p^2+1\) (Part~D4). Hence \(9\mid P(1)\) or
\(12\mid P(1)\) forces \(P=(X^2+p)^2\) if \(p\equiv2\pmod3\), and \(P=(X^2-p)^2\) if \(p\equiv1\pmod3\). Let
\(\pi\) be the Frobenius endomorphism of \(\widetilde A\). By the Cayley--Hamilton theorem \((\pi^2\pm p)^2=0\), and
\(\Q[\pi]\) is a product of fields by (I2)(e), so \(\pi^2=\mp p\). Then
\(\widetilde A(\F_p)=\ker(1-\pi)\) is killed by \((1+\pi)(1-\pi)=1\pm p\), that is, by \(p+1\), respectively
\(p-1\). As \(A(\Q)[m]\) injects into \(\widetilde A(\F_p)\) for \(m\) prime to \(p\) (I2)(b), a point of order
\(9\) needs \(p\equiv\mp1\pmod9\), and a point of order \(12\) needs \(p\equiv\mp1\pmod{12}\) (Part~D5). Choose \(p\)
inert in \(M\) with \(p\not\equiv\pm1\pmod9\) and \(p\equiv5,7\pmod{12}\): for \(M=\Q(\sqrt{-3})\) take \(p=5\), and
for \(M=\Q(i)\) take \(p=7\) (Part~D6); otherwise \(M\cap\Q(\zeta_{36})=\Q\), because the imaginary quadratic
subfields of \(\Q(\zeta_{36})\) are \(\Q(i)\) and \(\Q(\sqrt{-3})\), and Chebotarev's theorem in \(M(\zeta_{36})\)
gives infinitely many primes \(p\equiv5\pmod{36}\) inert in \(M\). Finally, for \(p\equiv5,7\pmod{12}\) no value
\(P(1)\) above is divisible by \(72\) (Part~D3), and a subgroup of order \(72\) of \(A(\Q)_{\mathrm{tors}}\) would
inject into \(\widetilde A(\F_p)\). \qed

\section{Nonisogenous factors}\label{typeN}

Let \(A/\Q\) be an abelian surface with a principal polarization \(\lambda\) defined over \(\Q\), simple over
\(\Q\) and of type (N), and let \(E\), \(E^\sigma\subset A_K\), \(n\), \(\psi_n\), \(\varphi_n\), \(\pi_E\) and
\(\Phi_n\) be as in (I3) and (I6). The argument of \cite[Proposition~2, step~1]{P2214} gives an exact sequence
\[0\to S\to A(\Q)_{\mathrm{tors}}\xrightarrow{\ \pi_E\ }E(K)_{\mathrm{tors}},\qquad S=A(\Q)_{\mathrm{tors}}\cap E:\]
if \(P\in A(\Q)\) and \(\pi_E(P)=0\), then \(P\in E^\sigma\), hence \(P=\tilde\sigma(P)\in\tilde\sigma(E^\sigma)=E\).
Every \(P_0\in S\) also lies in \(\tilde\sigma(E)=E^\sigma\), hence in \(E\cap E^\sigma=E[n]\); let \(x\in E(K)\) be \(P_0\) regarded as a point of \(E\), so that
\(P_0=\varphi_n(x,0)\). As \(\varphi_n\) is defined over \(\Q\), (I5) gives
\(P_0=\tilde\sigma(P_0)=\varphi_n(\tilde\sigma(x,0))=\varphi_n(0,\sigma x)\), so \((x,-\sigma x)\) lies in the kernel
of \(\varphi_n\), the graph of \(\psi_n\):
\[\psi_n(x)=-\sigma(x)\qquad\text{for every }x\in E(K)\text{ corresponding to a point of }S\tag{1}\]
(this is equation (1) in \cite[proof of Lemma~3]{P45}).

\begin{statement}{Theorem 5 (nonisogenous factors).}\label{thm-N}
In this situation, suppose that \((\Z/3)^2\subseteq A(\Q)\). Then \(K=\Q(\sqrt{-3})\), \(E[3]\subseteq E(K)\),
\(E(K)_{\mathrm{tors}}\cong\Z/3\times\Z/3\) or \(\Z/3\times\Z/6\), and \(\#A(\Q)[2^\infty]\le4\). If moreover
\(v_2(n)\le1\), then \(A(\Q)[2^\infty]\) is killed by \(2\).
\end{statement}

\emph{Proof.} Let \(W\cong(\Z/3)^2\) be a subgroup of \(A(\Q)\), and suppose that \(E(K)[3]\) is cyclic. The image
\(\pi_E(W)\cong W/(W\cap S)\) lies in the cyclic group \(E(K)[3]\), of order at most \(3\), so \(W\cap S\neq0\); as \(S\subseteq E[n]\), this gives \(3\mid n\).
Since \(W\cap S\) embeds in the cyclic group \(E(K)[3]\), we get \(W\cap S=\langle P_0\rangle\) of order \(3\) and
\(E(K)[3]=\langle P\rangle\) of order \(3\), with \(P\) corresponding to \(P_0\); by (1), \(\psi_n(P)=-\sigma(P)\).
The image \(\pi_E(W)\) is nonzero, so there is \(R\in W\) with \(\Phi_n(R)=(P',\sigma P')\in(\Res_{K/\Q}E)(\Q)\) and
\(0\neq P'\in\langle P\rangle\). Then \(\varphi_n(P',\sigma P')=nR=0\), so \(\psi_n(P')=\sigma(P')\), while
\(\psi_n(P')=-\sigma(P')\); hence \(2\sigma(P')=0\), a contradiction. So \(E(K)[3]\) is not cyclic, that is,
\(E[3]\subseteq E(K)\); by (I1), \(K=\Q(\sqrt{-3})\) and \(E(K)_{\mathrm{tors}}\cong\Z/3\times\Z/3\) or
\(\Z/3\times\Z/6\). The exact sequence, restricted to \(2\)-power torsion, gives
\(\#A(\Q)[2^\infty]\le\#E(K)[2^\infty]\cdot\#E(K)[2^\infty]\le4\).

Suppose now that \(v_2(n)\le1\) and \(E(K)[2^\infty]=\langle T\rangle\cong\Z/2\) (if \(E(K)[2^\infty]=0\), then
\(A(\Q)[2^\infty]=0\)). If \(v_2(n)=0\), the kernel of \(\varphi_n\) has odd order, and
\(A(\Q)[2^\infty]\cong(\Res_{K/\Q}E)(\Q)[2^\infty]=E(K)[2^\infty]\). If \(v_2(n)=1\), let \(\Gamma_2\) be the
\(2\)-primary part of \(\ker\varphi_n\), the graph of \(\psi_2=\psi_n|_{E[2]}\), and
\(B=\Res_{K/\Q}E/\Gamma_2\), with quotient map \(\varphi\) and \(\Phi\colon B\to\Res_{K/\Q}E\) such that
\(\Phi\varphi=[2]\) and \(\varphi\Phi=[2]\). The isogeny \(B\to A\) induced by \(\varphi_n\) has odd degree, so
\(A(\Q)[2^\infty]\cong B(\Q)[2^\infty]\). For \(R\in B(\Q)[2^\infty]\),
\(2R=\varphi(\Phi(R))\in\{0,\varphi(T,\sigma T)\}\), and \(\varphi(T,\sigma T)=0\) since \(\psi_2(T)=\sigma T\), the
unique nonzero \(K\)-rational point of \(E^\sigma[2]\). So \(2R=0\). \qed

\section{Endomorphism field of degree two: an inertia argument}\label{inertia}

In this section \(A/\Q\) is an abelian surface of type (S), simple over \(\Q\), with \(\Gamma\cong C_2\), so that
\(L\) is a quadratic field; \(\sigma\) is the generator of \(\Gamma\), \(\chi_L\colon G_{\Q}\to\{\pm1\}\) is the
character of \(L\), and \(D=\End^0(A_{\bar\Q})=\End^0(A_L)\cong M_2(\Q)\). For a prime \(\ell\) we write
\(\rho_\ell\colon G_{\Q}\to\GL(V_\ell(A))\) and \(T_\ell=T_\ell(A)\).

\begin{statement}{Lemma 6 (real multiplication).}\label{lemma-rm}
The action of \(\sigma\) on \(D\) is \(\Ad(R)\) for an element \(R\in D^\times\) with \(R^2=r\in\Q^\times\), \(r\) not
a square, and \(\End^0_{\Q}(A)=D^\Gamma=\Q(R)\cong\Q(\sqrt r)\). If \(L\) is imaginary, then \(r>0\), so
\(\End^0_{\Q}(A)\) is a real quadratic field.
\end{statement}

\emph{Proof.} By Skolem--Noether \(\sigma\) acts as \(\Ad(R)\) for some \(R\in D^\times\), and \(\Ad(R)^2=1\) gives
\(R^2=r\in\Q^\times\). As \(\sigma\neq1\), \(R\notin\Q^\times\), so \(R\) has a quadratic minimal polynomial and its
centralizer in \(D\cong M_2(\Q)\) is \(\Q[R]\). This centralizer is \(D^{\Ad(R)}=D^\Gamma=\End^0_{\Q}(A)\). If \(r\)
were a square, \(\Q[R]\cong\Q\times\Q\) would contain a nontrivial idempotent, and \(A\) would not be simple over
\(\Q\). So \(r\) is not a square and \(\End^0_{\Q}(A)=\Q(R)\cong\Q(\sqrt r)\).

Let \(L\) be imaginary and let \(\tau\in G_{\Q}\) be a complex conjugation; then \(\tau\notin G_L\), so \(\tau\)
acts on \(D\) as \(\Ad(R)\). Put \(\Lambda=H_1(A(\mathbb C),\Q)\). As a \(D\)-module, \(\Lambda\cong Y\otimes W\)
with \(D=\End(Y)\otimes1\) and \(\dim Y=\dim W=2\). The antiholomorphic involution \(\tau\) of \(A(\mathbb C)\)
induces \(\tau_*\in\GL(\Lambda)\) with \(\tau_*\varphi\tau_*^{-1}=\tau(\varphi)=R\varphi R^{-1}\) for
\(\varphi\in D\), so \(R^{-1}\tau_*\) centralizes \(D\) and \(\tau_*=R\,(1\otimes w)\) with \(w\in\GL(W)\). From
\(\tau_*^2=1\) we get \(r\,(1\otimes w^2)=1\), that is, \(w^2=1/r\). The complex structure \(J\) of
\(\Lambda\otimes\mathbb R\) commutes with \(D\), since endomorphisms are holomorphic, so \(J=1\otimes j\) with
\(j\in\End(W\otimes\mathbb R)\) and \(j^2=-1\); and \(\tau_*J=-J\tau_*\) gives \(wj=-jw\). Thus \(1,j,w,jw\) are
linearly independent and span a copy of the quaternion algebra \((-1,1/r)_{\mathbb R}\) inside
\(\End(W\otimes\mathbb R)\cong M_2(\mathbb R)\); by dimension it is all of \(M_2(\mathbb R)\), so it is split, and
hence \(1/r>0\). \qed

\begin{statement}{Lemma 7 (cyclic isogenies and the valuation of j).}\label{lemma-tate}
Let \(K'\) be a finite extension of \(\Q_p\), \(v\) the valuation of \(\bar K'\) with \(v(p)=1\), and \(f\colon E\to
E'\) a cyclic isogeny of degree \(N\) of elliptic curves over \(K'\) with \(v(j(E))<0\). Then
\(v(j(E'))=(s/d)\,v(j(E))\) for positive integers \(s\), \(d\) with \(sd=N\). In particular, if
\(v(j(E'))=v(j(E))\), then \(N\) is a square.
\end{statement}

\emph{Proof.} By (I8) there is \(q\in K'\) with \(v(q)=-v(j(E))>0\) and \(j(E)=j(q)\). Fix \(Q_0\in\bar K'\) with
\(Q_0^N=q\) and a primitive \(N\)-th root of unity \(\xi\in\bar K'\), and put \(K''=K'(Q_0,\xi)\), a complete field.
A Laurent series in \(T\) with coefficients in \(\Z[\zeta_N]\) and finitely many negative terms converges at
\(T=Q_0\) in \(K''\), since \(v(Q_0)>0\), and evaluation at \(T=Q_0\), \(\zeta_N\mapsto\xi\), is a ring homomorphism
on these series. As polynomials in \(X\), both sides of the identity of (I9) have such series as coefficients,
since \(c(n)\in\Z\) and \(\Phi_N\in\Z[X,Y]\). Applied to that identity with \(q^{1/N}=T\), so that
\(q^{s/d}=T^{s^2}\) and \(\zeta_d^b=\zeta_N^{sb}\) when \(sd=N\), it gives
\[\Phi_N(X,j(E))=\prod_{s,b,d}\bigl(X-j(\xi^{sb}Q_0^{s^2})\bigr)\qquad\text{in }K''[X].\]
The dual isogeny \(\hat f\colon E'\to E\) is cyclic of degree \(N\), so \(\Phi_N(j(E'),j(E))=0\) by (I9), and
\(j(E')=j(\xi^{sb}Q_0^{s^2})\) for some \(s,b,d\). By (I8),
\(v(j(E'))=-s^2v(Q_0)=-(s^2/N)\,v(q)=(s/d)\,v(j(E))\). If \(v(j(E'))=v(j(E))\neq0\), then \(s=d\) and \(N=s^2\). \qed

\begin{statement}{Theorem 8 (an inertia element of order four).}\label{thm-inertia}
Let \(A/\Q\) be an abelian surface, simple over \(\Q\), of type (S) with \(\Gamma\cong C_2\). Suppose that
\(\End^0_{\Q}(A)\cong\Q(\sqrt a)\) is a real quadratic field (\(a>1\) squarefree; by \hyperref[lemma-rm]{Lemma~6}
this holds whenever \(L\) is imaginary), and that \(L\) is ramified at an odd prime \(p\). Then there is an element
\(h\) of an inertia group at \(p\) in \(G_{\Q}\) such that \(\rho_\ell(h)^2=-1\) for every prime \(\ell\neq p\).
Consequently
\[A(\Q)[2^\infty]\subseteq A[2^\infty]^{h}\cong(\Z/2)^2,\qquad A(\Q)[\ell^\infty]=0\ \text{ for every prime
}\ell\notin\{2,p\}.\]
In particular \(\#A(\Q)[2^\infty]\le4\), and \(A(\Q)\) has no point of order \(4\).
\end{statement}

There is no hypothesis on the polarization (Step~2 of the proof uses a polarization of \(A\) defined over \(\Q\), of
arbitrary degree, which always exists) or on the torsion of \(A\).

\emph{Proof.} Let \(R\) be as in \hyperref[lemma-rm]{Lemma~6}. As \(R\in D^\Gamma=\End^0_{\Q}(A)\), we may scale it
so that \(R\in\End_{\Q}(A)\) and \(R^2=n\in\Z_{>0}\); then \(F:=\End^0_{\Q}(A)=\Q(R)\), \(n\in a\,\Q^{\times2}\) is not
a square, and we write \(\gamma\) for the nontrivial automorphism of \(F\).

\emph{Step 1: a twisting endomorphism.} The elements of \(D\cong M_2(\Q)\) that anticommute with \(R\) form a
\(2\)-dimensional \(\Q\)-subspace: in a basis in which \(R\) is the companion matrix of \(X^2-n\), they are the
matrices \(S\) with
\[R=\begin{pmatrix}0&n\\1&0\end{pmatrix},\qquad S=\begin{pmatrix}x&-ny\\y&-x\end{pmatrix},\qquad
S^2=(x^2-ny^2)\cdot1 .\]
Each nonzero \(S\) is invertible, as \(n\) is not a square. Fix such an \(S\), scaled into \(\End(A_L)\). Then
\(\sigma(S)=RSR^{-1}=-S\), that is, \(g(S)=\chi_L(g)S\) for \(g\in G_{\Q}\), and on \(V_\ell(A)\)
\[\rho_\ell(g)\,S\,\rho_\ell(g)^{-1}=\chi_L(g)\,S\qquad(g\in G_{\Q}).\tag{2}\]

\emph{Step 2: the \(F_\ell\)-structure.} Put \(F_\ell=F\otimes\Q_\ell\). Since \(F\) acts on
\(H_1(A(\mathbb C),\Q)\), a \(2\)-dimensional \(F\)-vector space, \(V_\ell(A)\) is free of rank \(2\) over
\(F_\ell\), and \(G_{\Q}\) acts \(F_\ell\)-linearly. For \(g\in G_{\Q}\) let \(t_g=\tr_{F_\ell}\rho_\ell(g)\) and
\(\delta_g=\det_{F_\ell}\rho_\ell(g)\). As \(SR=-RS\), \(S\) is \(\gamma\)-semilinear, so (2) gives
\[\gamma(t_g)=\chi_L(g)\,t_g .\tag{3}\]
Moreover \(\delta_g=\chi_\ell(g)\) (Ribet's lemma \cite{Ri}; we include the proof). Let \(\lambda\) be a
polarization of \(A\) defined over \(\Q\) and \(e\) the associated Weil pairing on \(V_\ell(A)\). The Rosati
involution is a positive involution of the totally real field \(F\), hence trivial on \(F\), so
\(e(\alpha x,y)=e(x,\alpha y)\) for \(\alpha\in F\). As \(\Tr_{F_\ell/\Q_\ell}\) is a perfect pairing, there is a
unique \(F_\ell\)-bilinear \(\Psi\colon V_\ell\times V_\ell\to F_\ell(1)\) with
\(e=\Tr_{F_\ell/\Q_\ell}\circ\Psi\); it is alternating, nondegenerate on each factor of \(F_\ell\), and by uniqueness
\(\Psi(gx,gy)=\chi_\ell(g)\Psi(x,y)\). On a free module of rank \(2\) this says \(\delta_g=\chi_\ell(g)\).

\emph{Step 3: potentially good reduction at \(p\).} Let \(E\subset A_L\) be an elliptic curve; it is defined over
\(L\) by (I3). Put \(E^\sigma=\tilde\sigma(E)\). For an idempotent \(e_E\in D\) with image \(E\), the idempotent
\(\sigma(e_E)=\tilde\sigma e_E\tilde\sigma^{-1}\) has image \(\tilde\sigma(E)\) and equals \(Re_ER^{-1}\), whose image
is \(R(E)\); so \(E^\sigma=R(E)\). If \(E^\sigma=E\), then \(R|_E\) would be an endomorphism of \(E\) whose square is
the non-square integer \(n\), and \(E\) would have complex multiplication, contrary to type (S). So
\(E^\sigma\neq E\), and the addition map \(E\times E^\sigma\to A_L\) is an isogeny over \(L\). The restriction
\(R|_E\colon E\to E^\sigma\) is an isogeny over \(L\) of degree \(n\): indeed \(R|_{E^\sigma}\circ R|_E=[n]_E\) has
degree \(n^2\), and \(R|_{E^\sigma}=\tilde\sigma\circ R|_E\circ\tilde\sigma^{-1}\), as \(R\) is defined over \(\Q\),
has the same degree as \(R|_E\). The kernel of \(R|_E\) is isomorphic to \(\Z/m\times\Z/mN\) for some \(m,N\ge1\),
so \(R|_E=\mu\circ[m]\) with \(\mu\colon E\to E^\sigma\) a cyclic isogeny of degree \(N\), defined over \(L\) because
\([m]\) is surjective and defined over \(L\). As \(n=m^2N\) is not a square, neither is \(N\).

Let \(\mathfrak p\) be the prime of \(L\) above \(p\), which is unique because \(p\) ramifies in \(L\), and let
\(v\) be the valuation of the completion \(L_{\mathfrak p}=L\otimes\Q_p\) with \(v(p)=1\). Suppose that \(E\) does
not have potentially good reduction at \(\mathfrak p\), that is, that \(E\) has potentially multiplicative
reduction there. Then \(v(j(E))<0\) by (I8). The automorphism \(\sigma\otimes1\) of \(L_{\mathfrak p}\) over
\(\Q_p\) preserves \(v\), so \(v(j(E^\sigma))=v(\sigma(j(E)))=v(j(E))\). By \hyperref[lemma-tate]{Lemma~7}, applied
over \(K'=L_{\mathfrak p}\) to \(\mu\), \(v(j(E^\sigma))=(s/d)\,v(j(E))\) with \(sd=N\); so \(s=d\) and \(N=s^2\) is a
square, a contradiction. Hence \(E\) has potentially good reduction at \(\mathfrak p\), and so does \(E^\sigma\),
which is \(L\)-isogenous to \(E\); by (I2)(a), so do \(E\times E^\sigma\) and \(A\), which is isogenous to it over
\(L\).

\emph{Step 4: an inertia element of order \(4\).} Fix a place of \(\bar\Q\) above \(p\), and let \(I_p\subset
G_{\Q}\) be its inertia group. As \(L/\Q\) is ramified at \(p\), there is \(g\in I_p\) with \(\chi_L(g)=-1\). Let
\(\ell\neq p\). By Step~3 and (I2)(d), \(\rho_\ell(g)\) has finite order, and its characteristic polynomial \(P_g\)
on \(V_\ell(A)\) has coefficients in \(\Z\) that do not depend on \(\ell\). By Step~2, \(\delta_g=\chi_\ell(g)=1\),
and by (3), \(t_g=y_\ell\sqrt a\) with \(y_\ell\in\Q_\ell\). Hence
\[P_g=\Norm_{F_\ell/\Q_\ell}\bigl(X^2-t_gX+1\bigr)=(X^2-t_gX+1)(X^2+t_gX+1)=(X^2+1)^2-c\,X^2,\qquad c=a\,y_\ell^2\]
(Part~E1 of the certificate). So \(c\in\Z\) does not depend on \(\ell\), and \(c/a\) is a square in \(\Q_\ell\) for
every \(\ell\neq p\). If \(c\neq0\), this forces \(c/a\in\Q^{\times2}\), since a rational number that is not a
square is not a square in \(\Q_\ell\) for infinitely many primes \(\ell\). The roots of \(P_g\) are roots of unity
\(\xi\) with \(\xi+\xi^{-1}=\pm\sqrt c\), so \(0\le c\le4\). Since \(a>1\) is squarefree,
\(c\in\{0\}\cup(a\Z^2\cap[1,4])\), that is, \(c=0\), or \((a,c)=(2,2)\), or \((a,c)=(3,3)\) (Part~E2). Then
\(P_g=\Phi_4^2\), \(\Phi_8\) or \(\Phi_{12}\), and as \(\rho_\ell(g)\) is semisimple, \(\rho_\ell(g)^2=-1\),
\(\rho_\ell(g)^4=-1\) or \(\rho_\ell(g)^6=-1\), respectively (Part~E3). Put \(h=g\), \(g^2\) or \(g^3\),
respectively. Then \(h\in I_p\) and \(\rho_\ell(h)^2=-1\) for every \(\ell\neq p\), since the case does not depend
on \(\ell\).

\emph{Step 5: fixed points.} Rational points are fixed by \(h\). Let \(\ell=2\), which is different from \(p\).
Via \(h\), \(T_2\) is a module over \(\Z_2[X]/(X^2+1)=\Z_2[i]\), a discrete valuation ring with uniformizer \(i-1\)
and residue field \(\F_2\); being torsion-free of \(\Z_2\)-rank \(4\), it is free of rank \(2\). Hence
\(A[2^\infty]^h=\ker\bigl(i-1\ \text{on}\ (\Q_2(i)/\Z_2[i])^2\bigr)\cong(\Z_2[i]/(i-1))^2\cong(\Z/2)^2\), which is
contained in \(A[2]\) (Part~E4). For odd \(\ell\neq p\), \((\rho_\ell(h)-1)(\rho_\ell(h)+1)=-2\) is invertible on
\(T_\ell\), so \(A[\ell^\infty]^h=0\) (Part~E5). \qed

\section{Proof of the Theorem and consequences for the census}\label{proof}

\emph{Proof of the Theorem.} Let \(A\) be as in the Theorem, with principal polarization \(\lambda\) defined over
\(\Q\), and \((\Z/3)^2\subseteq A(\Q)\). As \(\lambda\) has degree \(1\), \hyperref[thm-full3]{Theorem~1} applies,
and by its part (d) there are four cases.
\begin{itemize}
\item Type (N): \hyperref[thm-N]{Theorem~5} gives \(\#A(\Q)[2^\infty]\le4\).
\item Type (CM): \(9\mid\#A(\Q)_{\mathrm{tors}}\), and \(72\nmid\#A(\Q)_{\mathrm{tors}}\) by
\hyperref[thm-cm]{Theorem~4}; so \(8\nmid\#A(\Q)[2^\infty]\), that is, \(\#A(\Q)[2^\infty]\le4\).
\item Type (S) with \(\Gamma\cong D_6\): \(\#A(\Q)[3^\infty]\ge9\), and \hyperref[prop-ord3]{Proposition~3}(c)
gives \(\#A(\Q)[2^\infty]\le4\).
\item Type (S) with \(\Gamma\cong C_2\): \(L=\Q(\sqrt{-3})\), which is imaginary and ramified at the odd prime
\(3\). By \hyperref[lemma-rm]{Lemma~6}, \(\End^0_{\Q}(A)\) is a real quadratic field, and
\hyperref[thm-inertia]{Theorem~8} with \(p=3\) gives \(A(\Q)[2^\infty]\subseteq(\Z/2)^2\). \qed
\end{itemize}

\medskip
\emph{Proof of the Corollary.} The Jacobian of a curve of genus \(2\) over \(\Q\) is an abelian surface with a
principal polarization defined over \(\Q\). The groups \([6,12]=\Z/6\times\Z/12\), \([3,24]=\Z/3\times\Z/24\) and
\([2,6,6]=\Z/2\times\Z/6\times\Z/6\) have order \(72\), contain \((\Z/3)^2\), and have \(2\)-primary parts
\(\Z/2\times\Z/4\), \(\Z/8\) and \((\Z/2)^3\), of order \(8\) (Part~F of the certificate). \qed

\medskip
On the census, \([6,12]\), \([3,24]\) and \([2,6,6]\) become impossible in the class ``geometrically split, simple
over \(\Q\)'' (before this proof the three cells were empty and unexcluded). All three groups have examples split
over \(\Q\) (the census curves 6.12.a, 3.24.a and 2.6.6.a), so their combined geometrically split cells keep the
grades of the class ``split over \(\Q\)'' (at present ``exact'' for all three).

\section{Remarks}\label{remarks}

\begin{enumerate}
\item\label{rem-more} \emph{More on the case \(\Gamma\cong C_2\).} Let \(A\) be simple over \(\Q\), of type (S)
with \(\Gamma\cong C_2\) and \(L=\Q(\sqrt{-3})\) (by \hyperref[thm-full3]{Theorem~1}(d) this is the case
\(\Gamma\cong C_2\) of the Theorem). By \hyperref[lemma-rm]{Lemma~6} and \hyperref[thm-inertia]{Theorem~8} with
\(p=3\), \(A(\Q)\) has no point of order \(4\) and no point of prime order \(\ell\ge5\), so
\(A(\Q)_{\mathrm{tors}}\subseteq(\Z/2)^2\times A(\Q)[3^\infty]\), with no hypothesis on the polarization. For
instance \([3,12]\) and \([36]\) cannot occur as \(A(\Q)_{\mathrm{tors}}\) in this case.

\item\label{rem-gamma} \emph{Other groups \(\Gamma\).} The proof of \hyperref[thm-inertia]{Theorem~8} uses
\(\Gamma\cong C_2\) only through an element \(g\) of an inertia group at an odd prime \(p\) whose image in \(\Gamma\)
is an involution \(\Ad(R)\) with \(R^2\in\Q_{>0}\setminus\Q^{\times2}\). If \(\Gamma\) is larger and such an
element exists, the same steps over the fixed field \(L'\subset L\) of the image of \(g\), over which
\(\End^0_{L'}(A)=\Q(R)\) is a real quadratic field, give \(\#A(\Q)[2^\infty]\le4\). For example, if
\(\Gamma\cong D_6\) and \(\sqrt{-3}\in L\), the elements of an inertia group at \(3\) that act nontrivially on
\(\sqrt{-3}\) map to involutions of \(\Gamma\), and the three involutions of \(\Gamma\) are conjugate, so this recovers
\hyperref[prop-ord3]{Proposition~3}(c) when their common class of \(R^2\) is not a square. The hypothesis
\(R^2\notin\Q^{\times2}\) is needed: in \(\PGL_2(\Q)\) the classes of
\(\left(\begin{smallmatrix}0&-1\\1&1\end{smallmatrix}\right)\) and
\(\left(\begin{smallmatrix}0&1\\1&0\end{smallmatrix}\right)\) generate a dihedral group of order \(6\) whose
involutions have trace-zero representatives with \(R^2=1\). We do not use this remark.

\item\label{rem-hyp} \emph{The hypotheses of \hyperref[thm-inertia]{Theorem~8}.} (i) Real multiplication is
needed. The census curve 2.2.2.2.d (LMFDB 1521.b1, table \texttt{g2c\_curves\_new}),
\(y^2+(x^2+x)y=3x^5+16x^4+17x^3+x^2-3x\), is simple over \(\Q\) and geometrically split, with geometric
endomorphism algebra \(M_2(\Q)\), Sato--Tate group \(E_2\) (component group \(C_2\); the component group of the Sato--Tate group is
\(\Gal(L/\Q)\) \cite{FKRS}) and an imaginary quadratic
endomorphism algebra over \(\Q\) \cite{LMFDB}: it is of type (S), with \(\Gamma\cong C_2\) and imaginary
multiplication. The discriminant of this model is \(3^6\cdot13^6\), and its rational torsion
subgroup is \((\Z/2)^4\), of order \(16\) (certificate, Part~V). By Silverberg's theorem \cite[Theorem~2.4]{Sb}
(the proof of \hyperref[thm-full3]{Theorem~1}(c) works for \(A[m]\) with any \(m\ge3\)), \(L\) lies in
\(\Q(A[3])\cap\Q(A[4])\), which is unramified outside \(3\) and \(13\) by (I2)(a), as \(A\) has good reduction at
every other prime; so \(L\) is \(\Q(\sqrt{-3})\), \(\Q(\sqrt{13})\) or \(\Q(\sqrt{-39})\), and in each case ramified
at an odd prime. At a prime of good reduction that splits in \(L\), Frobenius lies in \(G_L\); by steps~1 and~2 of
the proof of \cite[Lemma~3]{P2214} it acts on \(V_\ell(A)\cong V_\ell(E)\otimes H\) as \(B\otimes c\) with
\(c\in\Q^\times\), so its characteristic polynomial is the square of a polynomial in \(\Q_\ell[X]\), hence of one in
\(\Z[X]\). Part~V checks that at the primes \(5\le p<200\) other than \(13\) the \(L\)-polynomial is a square if
\(p\) splits in \(\Q(\sqrt{13})\), and equals \((1-pT^2)^2+3b_p^2T^2\) with \(b_p\ge0\) otherwise, with \(b_7=b_{11}=b_{19}=2\) and \(b_{37}=4\). As \(7\) splits in
\(\Q(\sqrt{-3})\) and \(11\) in \(\Q(\sqrt{-39})\), this gives \(L=\Q(\sqrt{13})\). (For imaginary multiplication
the argument of (I4), which needs only \(\Norm(R)>0\), gives \(\delta_g=\chi_\ell(g)\chi_L(g)\) in place of Step~2;
so \(\delta_g=-1\) in Step~4, and an element with \(\rho_\ell(h)^2=-1\) need not exist.) (ii) The prime \(p\) must be
odd: the LMFDB curve 25600.f.512000.1 (25600.s1 in \texttt{g2c\_curves\_new}) is simple over \(\Q\) and
geometrically split, with Sato--Tate group \(J(E_1)\), endomorphism algebra \(\Q(\sqrt2)\) over \(\Q\) and
endomorphism field \(\Q(i)\) \cite{LMFDB}, and its rational torsion subgroup is \((\Z/2)^3\) (Part~V). (iii) The
bound \((\Z/2)^2\) is attained: the LMFDB curve 3969.d.250047.1 (3969.i1 in \texttt{g2c\_curves\_new}) is simple
over \(\Q\) and geometrically split, with Sato--Tate group \(J(E_1)\), endomorphism algebra \(\Q(\sqrt3)\) over
\(\Q\) and endomorphism field \(\Q(\sqrt{-3})\) \cite{LMFDB}, and its rational torsion subgroup is
\(\Z/2\times\Z/6\) (Part~V). It is in the case \(\Gamma\cong C_2\), \(L=\Q(\sqrt{-3})\) of
\hyperref[rem-more]{Remark~1}, with a rational point of order \(3\) but not \((\Z/3)^2\).

\item\label{rem-lmfdb} \emph{Validation (not used in the proof).} In the LMFDB genus-\(2\) tables
\texttt{g2c\_curves} and \texttt{g2c\_curves\_new} \cite{LMFDB} (the second contains every curve of the first that
is simple over \(\Q\) and geometrically split, with the label of the first table as \texttt{old\_label}), the curves
that are simple over \(\Q\) and geometrically split and have \((\Z/3)^2\subseteq J(\Q)_{\mathrm{tors}}\) have
rational torsion \([3,3]\), \([3,6]\) or \([3,9]\); so none has more than two rational points of \(2\)-power order.
Their Sato--Tate groups are \(J(E_3)\), the case \(\Gamma\cong D_6\), and \(N(\mathrm{SU}(2)\times\mathrm{SU}(2))\),
type (N): all curves with torsion \([3,6]\) have \(J(E_3)\), and all with torsion \([3,9]\) have
\(N(\mathrm{SU}(2)\times\mathrm{SU}(2))\). Types (CM) and (S) with \(\Gamma\cong C_2\), which
\hyperref[thm-full3]{Theorem~1}(d) also allows, do not occur among them. The two such curves of
\texttt{g2c\_curves}, 26244.d.314928.1 and 26244.e.472392.1 (torsion \([3,3]\), Sato--Tate group \(J(E_3)\)), have
endomorphism fields 6.0.34992.1 and 6.0.177147.2 (table \texttt{g2c\_endomorphisms}), sextic fields whose only
quadratic subfield is \(\Q(\sqrt{-3})\), as \hyperref[thm-full3]{Theorem~1}(d) predicts.

\item\label{rem-where} \emph{Where the hypotheses are used.} The principal polarization defined over \(\Q\) is used
only in type (N) (\hyperref[thm-N]{Theorem~5}), through the complement \(E^\sigma\) of \(E\) and the descent of the
gluing (I6); \hyperref[thm-full3]{Theorem~1} needs a polarization over \(\Q\) of degree prime to \(3\), and
\hyperref[thm-inertia]{Theorem~8} a polarization over \(\Q\) of any degree. \hyperref[lemma-shapes]{Lemma~2},
\hyperref[prop-ord3]{Proposition~3}, \hyperref[thm-cm]{Theorem~4} and \hyperref[lemma-rm]{Lemma~6} use no
polarization. Simplicity over \(\Q\) gives \(\Gamma\neq C_1\), the two conjugate elliptic curves in type (N), and
the non-square \(r\) of \hyperref[lemma-rm]{Lemma~6}.
\end{enumerate}

\section{Computations}\label{computations}

The certificate
\href{https://genus-2-torsion.github.io/data/knowledge/sources/qsimple_612_checks.m}{\texttt{qsimple\_612\_checks.m}}
(Magma V2.30-1, about \(4\) seconds) consists of assertions and ends with \texttt{ALL\_612\_CHECKS\_PASSED}; a
failed assertion stops it before that line. Its log is
\href{https://genus-2-torsion.github.io/data/knowledge/sources/qsimple_612_checks.log}{\texttt{qsimple\_612\_checks.log}}.
SHA-256 of the script:
\texttt{75ed93ac9e13e456\allowbreak 1343d07ce7539e94\allowbreak f46d627d3d7b9f77\allowbreak cea8b21ae7ec597a};
of the log:
\texttt{c6c60835c24a0751\allowbreak fe023c01beadcee8\allowbreak 78826e2f50d322ce\allowbreak 86f74b88bab12ebc}.

Part~A: \hyperref[thm-full3]{Theorem~1}; A1: the matrices
\(\left(\begin{smallmatrix}I&X\\0&\pm I\end{smallmatrix}\right)\) that multiply the symplectic form by \(\pm1\) are
those with \(X\) symmetric; A2: the \(28\) conjugacy classes of subgroups \(G\) of the group \(M\) of order \(54\)
with an element of multiplier \(-1\), their orders \(2\cdot3^k\), the inversion, the unique subgroup of index \(2\),
and their quotients; A3: the abstract groups \(U\sd C_2\) with \(\#U\le27\), and the members of the list of (I3) of
this form. Part~B: the polynomials and point counts of \hyperref[lemma-shapes]{Lemma~2}, including \(n=4\), and the
square roots used for \(\gamma_g\). Part~C: the residues behind \hyperref[lemma-shapes]{Lemma~2}(i) and (ii).
Part~D: \hyperref[thm-cm]{Theorem~4}; D1: all supersingular Weil polynomials of abelian surfaces over \(\F_p\) for
\(p\le23\); D2: their values at \(1\) are the lists of the proof, linked to D1 by assertions; D3--D6: the
divisibility statements. Part~E: \hyperref[thm-inertia]{Theorem~8}, Steps~4 and~5. Part~F: the three groups of the
Corollary. Part~V: the checks of \hyperref[rem-hyp]{Remark~3} (validation, not used in the proof). Every value
quoted in the text from the certificate is asserted, not only printed. Parts~A--F are exact computations with
integers, polynomials and finite groups; Part~V uses Magma's \texttt{TorsionSubgroup} and
\texttt{LPolynomial} for genus-\(2\) curves. The LMFDB data of Remarks~3 and~4 were read on 15 November 2026.

\begin{thebibliography}{LMFDB}

\bibitem[De]{De}
M. Deuring, \emph{Die Typen der Multiplikatorenringe elliptischer Funktionenk\"orper}, Abh. Math. Sem. Hansischen
Univ. \textbf{14} (1941), 197--272.

\bibitem[FKRS]{FKRS}
F. Fit\'e, K. S. Kedlaya, V. Rotger and A. V. Sutherland, \emph{Sato--Tate distributions and Galois endomorphism
modules in genus~\(2\)}, Compos. Math. \textbf{148} (2012), 1390--1442.

\bibitem[Ka]{Ka}
S. Kamienny, \emph{Torsion points on elliptic curves and \(q\)-coefficients of modular forms}, Invent. Math.
\textbf{109} (1992), 221--229.
\href{https://doi.org/10.1007/BF01232025}{doi:10.1007/BF01232025}.

\bibitem[KM]{KM}
M. A. Kenku and F. Momose, \emph{Torsion points on elliptic curves defined over quadratic fields}, Nagoya Math. J.
\textbf{109} (1988), 125--149: Theorem~(0.1).
\href{https://doi.org/10.1017/S0027763000002816}{doi:10.1017/S0027763000002816}.

\bibitem[La]{La}
S. Lang, \emph{Elliptic Functions}, 2nd edition, Graduate Texts in Mathematics~112, Springer, New York, 1987:
Chapter~5, Sections~1--3 (formula (1) of Section~2, Theorem~5 and the remark following it), and Chapter~13,
Theorem~12.

\bibitem[LMFDB]{LMFDB}
The LMFDB Collaboration, \emph{The L-functions and modular forms database}, \url{https://www.lmfdb.org}, 2026: the
tables \texttt{g2c\_curves}, \texttt{g2c\_curves\_new} (the extended genus-\(2\) data of Booker and Sutherland) and
\texttt{g2c\_endomorphisms}, read through the mirror devmirror.lmfdb.xyz.

\bibitem[Na]{Na}
F. Najman, \emph{Complete classification of torsion of elliptic curves over quadratic cyclotomic fields}, J. Number
Theory \textbf{130} (2010), 1964--1968: Section~1 and Theorem~2(ii).
\href{https://doi.org/10.1016/j.jnt.2009.12.008}{doi:10.1016/j.jnt.2009.12.008}.

\bibitem[P40]{P40}
Census proof of 14 November 2026, \emph{Rational \(5\)-torsion and the \(2\)-primary torsion of geometrically split
abelian surfaces}, \url{https://genus-2-torsion.github.io/proofs/qsimple-40-obstruction.html}: Theorem.

\bibitem[P45]{P45}
Census proof of 14 November 2026, \emph{No rational point of order \(45\) on a geometrically split abelian surface
simple over \(\Q\)}, \url{https://genus-2-torsion.github.io/proofs/qsimple-45-obstruction.html}: (I7), the proof
of Proposition~A (case~(v)), and the proof of Lemma~3 with its equation~(1).

\bibitem[P63]{P63}
Census proof of 21 October 2026, \emph{An obstruction to rational \(63\)-torsion on geometrically split abelian
surfaces}, \url{https://genus-2-torsion.github.io/proofs/qsimple-63-obstruction.html}: Sections~1--5,
Lemma~3, Section~9.1 and Lemma~9.

\bibitem[P2210]{P2210}
Census proof of 11 November 2026, \emph{Rational \(5\)-torsion and the \(2\)-rank of geometrically split
Jacobians}, \url{https://genus-2-torsion.github.io/proofs/qsimple-2210-obstruction.html}: (I1), (I6) and the proof
of Proposition~A.

\bibitem[P2214]{P2214}
Census proof of 10 November 2026, \emph{Rational \(7\)-torsion and the \(2\)-rank of geometrically split abelian
surfaces}, \url{https://genus-2-torsion.github.io/proofs/qsimple-2214-obstruction.html}: Section~1 ((I2) and
(I3)), Proposition~2 (step~1) and Lemma~3.

\bibitem[Ri]{Ri}
K. A. Ribet, \emph{Galois action on division points of abelian varieties with real multiplications}, Amer. J.
Math. \textbf{98} (1976), 751--804.

\bibitem[Sb]{Sb}
A. Silverberg, \emph{Fields of definition for homomorphisms of abelian varieties}, J. Pure Appl. Algebra
\textbf{77} (1992), 253--262: Theorem~2.4.

\bibitem[Si1]{Si1}
J. H. Silverman, \emph{The Arithmetic of Elliptic Curves}, Graduate Texts in Mathematics~106, Springer, New York,
1986: Proposition~VII.5.5.

\bibitem[Si2]{Si2}
J. H. Silverman, \emph{Advanced Topics in the Arithmetic of Elliptic Curves}, Graduate Texts in Mathematics~151,
Springer, New York, 1994: Theorem~II.6.1, Theorem~V.3.1 and Lemma~V.5.1.

\bibitem[ST]{ST}
J.-P. Serre and J. Tate, \emph{Good reduction of abelian varieties}, Ann. of Math. (2) \textbf{88} (1968),
492--517: \S1 (Theorem~1, Corollary~2 and Lemma~2) and \S2 (Theorem~2).
\href{https://doi.org/10.2307/1970722}{doi:10.2307/1970722}.

\bibitem[Ta]{Ta}
J. Tate, \emph{Endomorphisms of abelian varieties over finite fields}, Invent. Math. \textbf{2} (1966), 134--144:
Theorem~2(a).

\end{thebibliography}

\end{document}
